\documentclass[a4paper,11pt]{article}
\setlength{\textheight}{23.30cm}
\setlength{\textwidth}{16.5cm}
\setlength{\oddsidemargin}{0.2cm}
\setlength{\evensidemargin}{0.2cm}
\setlength{\topmargin}{0cm}
\setlength{\parindent}{0.4cm}


\usepackage{bbm}
\usepackage{amsfonts}
\usepackage{graphics,color}
\usepackage{amsmath}
\usepackage{amssymb}
\usepackage{mathrsfs}
\usepackage{cite}
\usepackage{verbatim}
\usepackage{float}
\usepackage{graphicx}
\usepackage{amsthm}
\usepackage{textcomp}
\usepackage{subfig}
\usepackage{esint}
\usepackage{enumerate}
\usepackage{hyperref}
%\usepackage{showkeys}


\newcommand{\red}{\color{red}}
\newcommand{\blue}{\color{blue}}
\newcommand{\black}{\color{black}}

\numberwithin{equation}{section}
\mathchardef\emptyset="001F

\newtheorem{Theorem}{Theorem}[section]
\newtheorem{Definition}[Theorem]{Definition}
\newtheorem{Proposition}[Theorem]{Proposition}
\newtheorem{Corollary}[Theorem]{Corollary}
\newtheorem{Lemma}[Theorem]{Lemma}


\newcommand{\nada}[1]{}
\newcommand{\numberset}{\mathbb}
\newcommand{\C}{\numberset{C}}
\newcommand{\eps}{\varepsilon}
\newcommand{\grad}{\nabla}
\newcommand{\grado}{{\rm deg}}
\newcommand{\graphofmapvk}{V_k}
\newcommand{\graphofmapu}{U}
\newcommand{\homogeneousmap}{\varphi}
\newcommand{\indice}{j}
\newcommand{\intervallounitario}{I}
\newcommand{\latus}{L}
\newcommand{\lift}{\Phi}
\newcommand{\M}{\numberset{M}}% la massa
\newcommand{\mappa}{u}
\newcommand{\mappabuona}{v}
\newcommand{\mres}{\mathbin{\vrule height 1.6ex depth 0pt width
0.13ex\vrule height 0.13ex depth 0pt width 1.3ex}}
\newcommand{\N}{\numberset{N}} 
\newcommand{\Om}{\Omega} 
\newcommand{\R}{\numberset{R}} 
\newcommand{\Rn}{\numberset{R}^n} 
\newcommand{\Suno}{{\mathbb S}^1}
\newcommand{\Z}{\numberset{Z}}
\newcommand{\rom}{\mathrm}

\DeclareMathOperator*{\esssup}{ess\,sup}

\theoremstyle{definition}
\newtheorem{Remark}[Theorem]{Remark}
\newtheorem{Example}[Theorem]{Example}

 \title{Note on Sobolev regularity for $p$-harmonic maps
}
\author{
Simone Carano, Nikos Katzourakis and Roger Moser}

\begin{document}
\maketitle

\begin{abstract} 

\end{abstract}
{\bf 2020 MSC:}\\
{\bf Key words and phrases:} Vectorial Calculus of Variations in $L^\infty$; higher order problems; Absolute Minimisers, Maximum Principle.

%%%%%%%%%%%%%%%%%%%%%%%%%%%%%%%%%%%%%%%%%%%%%%%%%%%%%%%%%%%%%%%%%%%%%%%%%
\section{Introduction}\label{sec:introduction}
%%%%%%%%%%%%%%%%%%%%%%%%%%%%%%%%%%%%%%%%%%%%%%%%%%%%%%%%%%%%%%%%%%%%%%%%




\begin{Lemma}
Let $\Om\subset\R^n$ an open set and $p\in[2,\infty)$. Assume that $u:\Om\to\R^N$ is a $p$-harmonic map, i.e. $u\in W^{1,p}_{\mathrm{loc}}(\Om,\R^N)$ satisfying
$$
\mathrm{div} (|\rom Du|^{p-2}\rom Du)=0\qquad\mbox{in }\Om,
$$
in the sense of distributions. Then, we have
$$
|\rom Du|^p\in W^{1,2}_{\mathrm{loc}}(\Om).
$$
\end{Lemma}
\begin{proof}
We observe
$$
\rom D|\rom Du|^p=\rom D(|\rom Du|^\frac p2)^2=2|\rom Du|^\frac p2 \rom D|\rom Du|^\frac p2.
$$
Therefore, since $u\in C^{1,\alpha}_\mathrm{loc}(\Om,\R^N)$, it is enough to show that $|\rom Du|^\frac p2\in W^{1,2}_{\mathrm{loc}}(\Om)$.\\
To this purpose, we can assume $p>2$, since for $p=2$ the result is clearly true.\\
Fix $U\Subset\Om$ open and $\eps>0$, and consider the minimum problem
$$
\min_{W^{1,p}_u(U,\R^N)}E_\eps,
$$  
where $W^{1,p}_u(U,\R^N)=u+W^{1,p}_0(U,\R^N)$ and
$$
E_\eps(v):=\int_U\left(|\mathrm Dv|^2+\eps^2\right)^{\frac p2}\qquad\forall v\in W^{1,p}(U,\R^N).
$$
Since $E_\eps$ is strictly convex, there exists a unique minimiser $u_\eps\in W^{1,p}_u(U,\R^N)$. Moreover, by regularity theory \cite{U}, we have $u_\eps\in C^\infty(U,\R^N)$.\\
As $\eps\to0$, $(E_\eps)_\eps$ is a sequence of equi-coercive and weakly lower semicontinuous functionals on $W^{1,p}(U,\R^N)$. Further, for every $v\in W^{1,p}(U,\R^N)$, $(E_\eps(v))_\eps$ is monotone decreasing with limit $E(v):=\int_U|Dv|^p$. Therefore, from the theory of $\Gamma$-convergence (see \cite[Section 2.4]{Br}), we have that
 $E_\eps\stackrel\Gamma\to E$ w.r.t. the weak topology of $W^{1,p}(U,\R^N)$. As a consequence, since $u$ is the unique minimiser of $E$ in $W^{1,p}_u(U,\R^N)$, the full sequence $u_\eps$ converges to $u$ weakly in $W^{1,p}(U,\R^N)$. We have also the convergence of their corresponding norms, since
\begin{align*}
E(u)=\lim_\eps E_\eps(u_\eps)&=\lim_\eps\int_U\left(|\mathrm Du_\eps|^2+\eps^2\right)^{\frac p2}\\
&\geq\limsup_\eps\int_U|\mathrm Du_\eps|^p\\
&\geq\liminf_\eps\int_U|\mathrm Du_\eps|^p\\
&\geq\int_U|\mathrm Du|^p=E(u)
\end{align*}
hence, the convergence of $u_\eps$ to $u$ is actually strong, namely
\begin{align}\label{strong conv}
u_\eps\to u\quad\mbox{in }W^{1,p}(U,\R^N).
\end{align} 
Now, we aim at deriving some uniform energy estimates for $u_\eps$.
 In the following, the symbols $\cdot$ and $ \langle\cdot,\cdot\rangle$ stand for the inner products in $\R^n$ and $\R^N$ respectively, while : stands for the hilbert-Schmidt inner product in $\R^{N\times n}$.\\
We start by considering the Euler-Lagrange equation corresponding to $E_\eps$:
$$
\rom{div}\big(\left(|\rom Du_\eps|^2+\eps^2\right)^{\frac{p-2}2}\rom Du_\eps\big)=0\quad\mbox{in }U.
$$
Since $u_\eps\in C^\infty(U,\R^N)$, we can differentiate the previous equation with respect to $x_i$, obtaining
\begin{align*}
0&=\partial_i\left[\rom{div}\big(\left(|\rom Du_\eps|^2+\eps^2\right)^{\frac{p-2}2}\rom Du_\eps\big)\right]\\
&=\rom{div}\left((p-2)\left(|\rom Du_\eps|^2+\eps^2\right)^{\frac{p-4}2}(\rom Du_\eps:\rom D\partial_i u_\eps)\rom Du_\eps+\left(|\rom Du_\eps|^2+\eps^2\right)^{\frac{p-2}2}\rom D\partial_iu_\eps\right).
\end{align*}
By taking the inner product with $\partial_i u_\eps$, we have
\begin{equation}\label{inner}
\begin{aligned}
0&=\left\langle\rom{div}\left((p-2)\left(|\rom Du_\eps|^2+\eps^2\right)^{\frac{p-4}2}(\rom Du_\eps:\rom D\partial_i u_\eps)\rom Du_\eps+\left(|\rom Du_\eps|^2+\eps^2\right)^{\frac{p-2}2}\rom D\partial_iu_\eps\right),\partial_iu_\eps\right\rangle\\
&=\rom{div}\left((p-2)\left(|\rom Du_\eps|^2+\eps^2\right)^{\frac{p-4}2}(\rom Du_\eps:\rom D\partial_i u_\eps)\langle\rom Du_\eps,\partial_iu_\eps\rangle+\left(|\rom Du_\eps|^2+\eps^2\right)^{\frac{p-2}2}\langle\rom D\partial _iu_\eps,\partial_iu_\eps\rangle\right)-\\
&\quad-(p-2)\left(|\rom Du_\eps|^2+\eps^2\right)^{\frac{p-4}2}(\rom Du_\eps:\rom D\partial_i u_\eps)^2-\left(|\rom Du_\eps|^2+\eps^2\right)^{\frac{p-2}2}|\rom D\partial_iu_\eps|^2
\end{aligned}
\end{equation}
Thus, testing the previous equation with $\eta\in C^\infty_c(U)$, we get
\begin{equation}\label{test}
\begin{aligned}
&\int_U\eta^2\left((p-2)\left(|\rom Du_\eps|^2+\eps^2\right)^{\frac{p-4}2}(\rom Du_\eps:\rom D\partial_i u_\eps)^2+\left(|\rom Du_\eps|^2+\eps^2\right)^{\frac{p-2}2}|\rom D\partial_iu_\eps|^2\right)\\
=&-2\int_U\eta \rom D\eta\cdot\left((p-2)\left(|\rom Du_\eps|^2+\eps^2\right)^{\frac{p-4}2}(\rom Du_\eps:\rom D\partial_i u_\eps)\langle\rom Du_\eps,\partial_iu_\eps\rangle+\left(|\rom Du_\eps|^2+\eps^2\right)^{\frac{p-2}2}\langle\rom D\partial _iu_\eps,\partial_iu_\eps\rangle\right).
\end{aligned}
\end{equation}
Using the Cauchy-Schwartz inequality, the right hand side is bounded above by
\begin{align*}
&\frac12\int_U\eta^2(p-2)\left(|\rom Du_\eps|^2+\eps^2\right)^{\frac{p-4}2}(\rom Du_\eps:\rom D\partial_i u_\eps)^2+2(p-2)\int_U|\rom D\eta|^2\left(|\rom Du_\eps|^2+\eps^2\right)^{\frac{p-4}2}|\rom Du_\eps|^2|\partial_iu_\eps|^2+\\
&\quad+\frac12\int_U\eta^2\left(|\rom Du_\eps|^2+\eps^2\right)^{\frac{p-2}2}|\rom D\partial_i u_\eps|^2+2\int_U|\rom D\eta|^2\left(|\rom Du_\eps|^2+\eps^2\right)^{\frac{p-2}2}|\partial_iu_\eps|^2,
\end{align*}
which in turn, using $|\rom Du_\eps|^2\leq\left(|\rom Du_\eps|^2+\eps^2\right)^{\frac12}$, is bounded by
\begin{align*}
&\frac12\int_U\eta^2\left((p-2)\left(|\rom Du_\eps|^2+\eps^2\right)^{\frac{p-4}2}(\rom Du_\eps:\rom D\partial_i u_\eps)^2+\left(|\rom Du_\eps|^2+\eps^2\right)^{\frac{p-2}2}|\rom D\partial_i u_\eps|^2\right)+\\
&\quad+2(p-1)\int_U|\rom D\eta|^2\left(|\rom Du_\eps|^2+\eps^2\right)^{\frac{p-2}2}|\partial_iu_\eps|^2.
\end{align*}
This last quantity provides an estimates from above for the left hand side of \eqref{test}, so that
\begin{align*}
&\frac12\int_U\eta^2\left((p-2)\left(|\rom Du_\eps|^2+\eps^2\right)^{\frac{p-4}2}(\rom Du_\eps:\rom D\partial_i u_\eps)^2+\left(|\rom Du_\eps|^2+\eps^2\right)^{\frac{p-2}2}|\rom D\partial_i u_\eps|^2\right)\\
\leq\,&2(p-1)\int_U|\rom D\eta|^2\left(|\rom Du_\eps|^2+\eps^2\right)^{\frac{p-2}2}|\partial_iu_\eps|^2.
\end{align*}
In particular,
\begin{align*}
\int_U\eta^2\left(|\rom Du_\eps|^2+\eps^2\right)^{\frac{p-4}2}(\rom Du_\eps:\rom D\partial_i u_\eps)^2
\leq4\,\frac{p-1}{p-2}\int_U|\rom D\eta|^2\left(|\rom Du_\eps|^2+\eps^2\right)^{\frac{p-2}2}|\partial_iu_\eps|^2.
\end{align*}
Now, we observe
$$
\partial_i \left(|\rom Du_\eps|^2+\eps^2\right)^{\frac p4}=p\left(|\rom Du_\eps|^2+\eps^2\right)^{\frac{p-4}4}\rom Du_\eps:\rom D\partial_iu_\eps,
$$
so that
$$
\int_U\eta^2\left(\partial_i \left(|\rom Du_\eps|^2+\eps^2\right)^{\frac p4}\right)^2\leq4p^2\,\frac{p-1}{p-2}\int_U|\rom D\eta|^2\left(|\rom Du_\eps|^2+\eps^2\right)^{\frac{p-2}2}|\partial_iu_\eps|^2.
$$
For any open subset $U'\Subset U$, we can choose $\eta\in C^\infty_c(U)$ such that $\eta=1$ on $U'$. Hence, by the H\"older inequality, setting $C(p,U,U'):=4p^2\,\frac{p-1}{p-2}\|\rom D\eta\|^2_{L^\infty(U)}$, we obtain
\begin{align}\label{deriv est}
\int_{U'}\left(\partial_i \left(|\rom Du_\eps|^2+\eps^2\right)^{\frac p4}\right)^2\leq C(p,U,U')E_\eps(u_\eps)^\frac{p-2}{p}\|\partial_iu_\eps\|^2_{L^p(U)}.
\end{align}
Thanks to \eqref{strong conv}, we have $\left(|\rom Du_\eps|^2+\eps^2\right)^{\frac p4}\to |\rom Du|^{\frac{p}{2}}$ in $L^2(U)$. Furthermore, by recalling that $E_\eps(u_\eps)\to E(u)=\|\rom Du\|^p_{L^p(U)}$ as $\eps\to0$, we can pass to the limit in \eqref{deriv est} and conclude
$$
\left\|\partial_i|\rom Du|^\frac p2\right\|^2_{L^2(U')}\leq\liminf_{\eps\to0}\int_{U'}\left(\partial_i \left(|\rom Du_\eps|^2+\eps^2\right)^{\frac p4}\right)^2\leq C(p,U,U')\|\rom Du\|^{p}_{L^p(U)}.
$$
Therefore, since $U'$ and $U$ were arbitrary, we have just shown that $|\rom Du|^\frac p2\in W^{1,2}_\mathrm{loc}(\Om)$.\\
\end{proof}



%%%%%%%%%%%%%        BIBLIOGRAFIA        %%%%%%%%%%%%%%%%%%%%%%%%%%%%
\begin{thebibliography}{AA} 
\addcontentsline{toc}{chapter}{Bibliografia} 

\bibitem{A}
G. Aronsson, {\it Minimization problems for the functional $\sup_xF(x, f(x), f'(x))$ I, II, III}, Arkiv
fur Mat. 33 - 53 (1965); 409 - 431 (1966); 509 - 512 (1969).

\bibitem{BBDS}
{A. Balci, L. Behn, L. Diening, J. Storn,} {\it {Examples of \({p}\)-harmonic Maps}}, {SIAM Journal on Mathematical Analysis} {\bf 58}(1), 260-275 (2026).

\bibitem{BJW}
E.N. Barron, R.R. Jensen, C.Y. Wang, {\it The Euler equation and absolute minimizers of $L^\infty$ functionals}, Arch. Rational
Mech. Anal. {\bf 157}, 255-283 (2001).

\bibitem{Br}
A. Braides, {\it A handbook of $\Gamma$-convergence},
handbook of Differential Equations: Stationary Partial Differential Equations,
North-holland,
{\bf 3}, 101-213 (2006).

\bibitem{BDP}
 C. Brizzi, L. De Pascale, \emph{A property of Absolute Minimizers in $ L^{\infty}$  Calculus of Variations and of solutions of the Aronsson-Euler equation}, Advances in Differential Equations {\bf 28} (3/4), 287-310 (2023).

\bibitem{CKM}
S. Carano, N. Katzourakis, R. Moser, \emph{
Existence, uniqueness and characterisation of vector-valued absolute minimisers for a second order $L^\infty$-variational problem}, Preprint arXiv https://arxiv.org/abs/2504.04181 (2025).

\bibitem{CM}
S. Carano, R. Moser, \emph{
An $L^\infty$-variational problem involving the Fractional Laplacian}, Preprint arXiv https://arxiv.org/abs/2510.14476 (2025).


%\bibitem{Dac}
%B. Dacorogna, {\it Direct Methods in the Calculus of Variations}, 2nd ed., Appl. Math. Sci. 78, Springer, New York, 2008.

%\bibitem{D}
%J. M. Danskin, {\it The theory of max-min with applications}, J. SIAM Appl. Math. {\bf 14}(4), 641-665 (1966).

\bibitem{DK}
B. Dutton, N. Katzourakis, {\it On Second-Order $L^\infty$ Variational Problems with Lower-Order Terms}, Discrete Contin. Dyn. Syst., {\bf 49}, 1-21 (2026).

\bibitem{GT}
D. Gilbarg, N. Trudinger, {\it Elliptic Partial Differential Equations of Second Order}.
Classics in Mathematics, reprint of the 1998 edition, Springer.

%\bibitem{GM}
%M. Giaquinta, L. Martinazzi, {\it An Introduction to the Regularity Theory for Elliptic Systems, harmonic Maps and Minimal Graphs}, Publications of the Scuola Normale Superiore (2018).

\bibitem{K1}
N. Katzourakis, {\it $L^\infty$ Variational Problems for maps and the Aronsson PDE System}, J. Differential
Equations {\bf 253}(7), 2123 - 2139 (2012).

\bibitem{KSIAM}
N. Katzourakis, {\it An $L^\infty$ regularisation strategy to the inverse source identification problem for elliptic equations}, SIAM Journal Math, Analysis {\bf 51}(2), 1349-1370 (2019).

%\bibitem{K2}
%N. Katzourakis. {\it Explicit 2D $\infty$-harmonic maps whose interfaces have junctions
%and corners}, C. R. Math. Acad. Sci. Paris {\bf 351}, 677–680 (2013).

\bibitem{Kbook}
N. Katzourakis, {\it An Introduction to Viscosity Solutions for Fully Nonlinear PDE with Applications
to Calculus of Variations in $L^\infty$}, Springer Briefs in Mathematics,  DOI 10.1007/978-
3-319-12829-0 (2015).

\bibitem{K3}
N. Katzourakis, {\it $\infty$-minimal submanifolds}, Proc. Amer. Math. Soc. {\bf 142}, 2797-2811 (2014).

\bibitem{K4}
N. Katzourakis, {\it On the structure of $\infty$-harmonic maps}, Comm. Partial Differential
Equations {\bf 39}, 2091–2124 (2014).

%\bibitem{K5}
%N. Katzourakis. {\it Nonuniqueness in vector-valued calculus of variations in $L^\infty$ and some linear elliptic systems}, Commun. Pure Appl. Anal. {\bf 14}, 313–327 (2015).

%\bibitem{K6}
%N. Katzourakis. {\it A characterisation of $\infty$-harmonic and $p$-harmonic maps via affine variations in $L^\infty$}, Electron. J. Differential Equations , {\bf 29} (2017).

\bibitem{KM}
N. Katzourakis, R. Moser, {\it Existence, uniqueness and structure of second order absolute minimisers}, ARMA {\bf 231}(3), 1615-1634 (2019).

\bibitem{KM1}
N. Katzourakis, R. Moser, {\it Minimisers of supremal functionals and mass-minimising 1-currents}, Calc. Var.{ \bf 64}(26) (2025).

\bibitem{KM2}
N. Katzourakis, R. Moser, {\it Variational problems in $L^\infty$ involving semilinear second order differential operators}, ESAIM: COCV, {\bf 29}(76) (2023).

\bibitem{KM3}
N. Katzourakis, R. Moser. {\it Existence, uniqueness and characterisation of local minimisers in higher order Calculus of Variations in $L^\infty$}, to appear in Analysis \& PDE (2026).

\bibitem{KP1}
N. Katzourakis, T. Pryer, {\it On the numerical approximation of $p$-Biharmonic and $\infty$-Biharmonic functions},
Numerical Methods for Partial Differential Equations, {\bf 35}(1), 155-180 (2018).


\bibitem{KP2}
N. Katzourakis, T. Pryer, {\it 2nd order $L^\infty$ Variational Problems and the $\infty$-Polylaplacian},
Advances in Calculus of Variations {\bf 13}, 115-140 (2020).



\bibitem{MWZ}
 Q. Miao, C. Wang, Y. Zhou, \emph{Uniqueness of Absolute Minimizers for $ L^{\infty}$-Functionals Involving hamiltonians $h(x,p)$}, Archive for Rational Mechanics and Analysis \textbf{223} (1), 141--198 (2017).

\bibitem{PWZ} F. Peng, C. Wang, Y. Zhou, \emph{Regularity of absolute minimizers for continuous convex hamiltonians}, Journal of Differential Equations, Volume 274 (2021), 1115-1164.

% \bibitem{PZ} F. Prinari, E. Zappale, \emph{A Relaxation Result in the Vectorial Setting and Power Law Approximation for Supremal Functionals}, J Optim. Theory Appl., \textbf{186}, 412--452 (2020).

%\bibitem{RZ} A. M. Ribeiro, E. Zappale, \emph{Existence of minimisers for nonlevel convex functionals}, SIAM J. Control Opt., \textbf{52} (5), 3341--3370 (2014).

\bibitem{U}
K. Uhlenbeck, {\it Regularity for a class of non-linear elliptic systems}, Acta Math. {\bf138}, 219--240 (1977).



\end{thebibliography}

\vspace{3mm}
\textsc{Simone Carano}\\
Department of Mathematics and Statistics, University of Reading\\
Whiteknights Campus, Pepper Lane, Reading RG6 6AX, UK\\
E-mail: s.carano@reading.ac.uk\\

\textsc{Nikos Katzourakis}\\
Department of Mathematics and Statistics, University of Reading\\
Whiteknights Campus, Pepper Lane, Reading RG6 6AX, UK\\
E-mail: n.katzourakis@reading.ac.uk\\

\textsc{Roger Moser}\\
Department of Mathematical Sciences, University of Bath\\
Bath BA2 7AY, UK\\
E-mail: r.moser@bath.ac.uk





\end{document}











